\documentclass[10pt,a4paper]{amsart}
\usepackage[margin=19mm]{geometry}
\usepackage{amsmath,amssymb,mathtools}
\usepackage[scr=rsfs,cal=euler]{mathalfa}
\usepackage{xfrac}
\usepackage[unicode,hidelinks,bookmarks=false]{hyperref}
\newcommand{\ie}{i.e.\ }
\newcommand{\cf}{cf.\ }
\newcommand{\resp}{respectively\ }
\newcommand{\Subsection}{\subsection}
\providecommand{\nobreakdash}{\nobreak\hbox{-}}
\setlength{\emergencystretch}{2em}
\multlinegap0pt
\usepackage{tikz}
\usepackage{amssymb}
\usepackage{algorithm}\usepackage{algpseudocode}
\usepackage{xcolor}
\usepackage[normalem]{ulem}
\newtheorem{proposition}{Proposition}
\newcommand{\R}{\mathbb{R}}
\title{Wasserstein-Based Identification of Metastable States in time series Data via Change Point Detection and Segment Clustering}
\author{David Gentile, Joshua Huang, James M. Murphy}
\date{Source: arXiv:2603.01989v1 (CC BY 4.0)}
\begin{document}
\maketitle

\noindent\emph{Source and licence.} Wasserstein-Based Identification of Metastable States in time series Data via Change Point Detection and Segment Clustering. Version arXiv:2603.01989v1 (CC BY 4.0), distributed by arXiv under the Creative Commons Attribution 4.0 International licence, http://creativecommons.org/licenses/by/4.0/, as recorded in the arXiv licence metadata for this exact version and shown on the abstract page https://arxiv.org/abs/2603.01989v1. Full licence notice: \texttt{licenses/arxiv-2603.01989-CC-BY-4.0.txt}.\par\smallskip

\noindent\textit{Editorial scope.} Counted unit: the article's proposition proposition (source0.tex lines 448--451). The article's complete original proof of it is delivered with it (source0.tex lines 452--454, one \texttt{proof} environment of 3 lines). No mathematics is written, repaired or re-derived: the statement and the proof are the article's own lines, in the article's own notation, and no retained line is substituted or re-flowed.  No further source-local context is needed: every cross-reference of every retained line resolves inside the two delivered spans.  Nothing was omitted from any retained span, so the package carries no omission record.  No retained line contains a citation, so the wrapper carries no bibliography and no bibliography entry.  No retained line contains a figure environment, an image-inclusion command or drawing code, so no image is delivered and the package declares no asset dependency.  Editorial formatting only, in addition: an AMS \texttt{amsart} preamble that restates the article's own declarations and macros used by the retained lines, namely the environments \texttt{proposition} on separate counters, together with the standard packages those lines need.  The retained environments print in this document as Proposition 1 in order of appearance, whereas the article prints the same items with the numbering of its own full document.  The complete native source of the article as distributed in the arXiv e-print for this exact version is bundled unchanged as \texttt{sources/195586/source0.tex}.
\begin{proposition}
For fixed \(w\) and \(q\), for inputs \(\{\vec{X_{t}}\}_{t=1}^T \subset \R^D\), the runtime complexity of Algorithm 1 is \(\mathcal{O}(TD).\)
{Fix \(w \in \{1, \dots \lfloor T/2\rfloor\}, q \in (0,1)\)}, and assume that the number of segments \(N\) scales proportionally to the length of the input data \(X = \{X_t\}_{t=1}^T\).
\end{proposition}

\begin{proof}
    The claim follows directly from the construction of the algorithm: for each \(t = w, \dots, T - w\), we perform two sorting operations, which are \(\mathcal{O}(w\log w)\), but \(\mathcal{O}(1)\) with respect to the size and dimension of the input data. Likewise, the subroutine identifying the inflection points of the identified segments is no more complex than sorting the empirical distribution of each segment. Since the number of segments is assumed to grow linearly in the number of data points, we amortize by considering the expected length of each segment, which will be \(T/N\). Performing \(N \propto T\) sorting jobs on a list of size \(T/N\) has complexity \(\mathcal{O}(N\frac{T}{N}\log(\frac{T}{N}) = \mathcal{O}(T)\). Because this procedure is carried out over each of the \(D\) dimensions of the data, the algorithm has total asymptotic time complexity \(\mathcal{O}(TD).\)
\end{proof}

\end{document}
